% ============================================================
% Paper E: Formal Definitions and Theorems
% ============================================================

\section{Preliminaries and Problem Definition}\label{sec:formalization}

\subsection{Notation}

Let $\mathcal{I} = \{i_1, i_2, \ldots, i_m\}$ be a finite set of items and
$D = [d_1, d_2, \ldots, d_N]$ be a temporally ordered sequence of transactions,
where each $d_t \subseteq \mathcal{I}$.
A \emph{pattern} (itemset) is any non-empty $P \subseteq \mathcal{I}$.

\subsection{Support Time Series and Dense Intervals}

\begin{definition}[Support Time Series]\label{def:support_ts}
For pattern $P$, window size $W \in \mathbb{Z}_{>0}$, and transaction sequence $D$ of length $N$,
the \emph{support time series} $s_P^W = [s_P^W(0), \ldots, s_P^W(N-W)]$ is defined as
\begin{equation}
s_P^W(t) = |\{d_j \in D \mid t \leq j < t+W,\; P \subseteq d_j\}|.
\end{equation}
\end{definition}

\begin{definition}[Dense Interval]\label{def:dense_interval}
Given threshold $\theta \in \mathbb{Z}_{>0}$, a \emph{dense interval} of pattern $P$ at window size $W$
is a maximal contiguous range $[a, b] \subseteq [0, N-W]$ such that
$s_P^W(t) \geq \theta$ for all $t \in [a, b]$.
Let $\mathcal{D}(P, W, \theta)$ denote the set of all dense intervals.
\end{definition}

\subsection{Dense Coverage Score}

\begin{definition}[Dense Coverage Score]\label{def:coverage_score}
For pattern $P$ at window size $W$ with threshold $\theta$,
the \emph{Dense Coverage Score} is
\begin{equation}
\mathrm{DCS}(P, W, \theta) = \sum_{[a,b] \in \mathcal{D}(P, W, \theta)}
  \left( (b - a + 1) \cdot \bar{s}_P^W([a,b]) \right),
\end{equation}
where $\bar{s}_P^W([a,b]) = \frac{1}{b-a+1} \sum_{t=a}^{b} s_P^W(t)$
is the mean support within the interval.
\end{definition}

The DCS jointly rewards \emph{temporal extent} (long intervals) and
\emph{density magnitude} (high support values), capturing
patterns that are both persistently and strongly dense.

\begin{remark}
Equivalently, $\mathrm{DCS}(P, W, \theta) = \sum_{[a,b]} \sum_{t=a}^{b} s_P^W(t)$,
i.e., the total accumulated support across all dense intervals.
This formulation avoids the threshold sensitivity of simple interval counting.
\end{remark}

\subsection{Dyadic Scale Hierarchy}

\begin{definition}[Dyadic Scale Hierarchy]\label{def:dyadic}
Given a base window size $W_0 \in \mathbb{Z}_{>0}$ and maximum level
$L = \lfloor \log_2(N / W_0) \rfloor$, the \emph{Dyadic Scale Hierarchy}
is the sequence of window sizes
\begin{equation}
\mathcal{W} = \{W_\ell = 2^\ell \cdot W_0 \mid \ell = 0, 1, \ldots, L\}.
\end{equation}
At each level $\ell$, we set the threshold proportionally:
$\theta_\ell = \lceil \theta_0 \cdot 2^\ell \rceil$,
where $\theta_0$ is the base threshold.
This maintains a constant \emph{density ratio} $\theta_\ell / W_\ell \approx \theta_0 / W_0$.
\end{definition}

\begin{definition}[Multi-Scale Dense Coverage Score]\label{def:ms_dcs}
The \emph{Multi-Scale Dense Coverage Score} aggregates across the dyadic hierarchy:
\begin{equation}
\mathrm{MSDCS}(P) = \sum_{\ell=0}^{L} \omega_\ell \cdot \mathrm{DCS}(P, W_\ell, \theta_\ell),
\end{equation}
where $\omega_\ell = 1 / (L+1)$ (uniform) or $\omega_\ell \propto 1/\ell$
(fine-scale preference).
\end{definition}

\subsection{Scale-Space Dense Ridge}

\begin{definition}[Scale-Space Dense Ridge]\label{def:ridge}
Consider the \emph{dense indicator function}
$\phi_P(\ell, t) = \mathbb{1}[s_P^{W_\ell}(t) \geq \theta_\ell]$
defined over the scale-position plane $\{0,\ldots,L\} \times \{0,\ldots,N-W_L\}$.
A \emph{Scale-Space Dense Ridge} for pattern $P$ is a connected component
$R \subseteq \{(\ell, t) : \phi_P(\ell, t) = 1\}$ that spans at least
$\ell_{\min}$ consecutive scale levels, i.e.,
\begin{equation}
\max_{(\ell,t) \in R} \ell - \min_{(\ell,t) \in R} \ell \geq \ell_{\min} - 1.
\end{equation}
The \emph{ridge strength} is $\rho(R) = \sum_{(\ell,t) \in R} s_P^{W_\ell}(t)$.
\end{definition}

Ridges identify temporal regions where density persists across multiple
scales, providing scale-invariant evidence of genuine dense behavior
(as opposed to noise-induced density at a single scale).

\subsection{Branch-and-Bound Dense Pruning}

\begin{definition}[Dense Coverage Upper Bound]\label{def:upper_bound}
For a candidate pattern $P$ with current itemset $C \subseteq P$,
define the \emph{upper bound} on DCS as
\begin{equation}
\overline{\mathrm{DCS}}(P, W, \theta) \leq \mathrm{DCS}(C, W, \theta),
\end{equation}
since adding items to $C$ can only reduce support:
$s_P^W(t) \leq s_C^W(t)$ for all $t$ when $C \subseteq P$.
\end{definition}

\subsection{Top-k Dense Pattern Mining Problem}

\begin{definition}[Top-k Dense Pattern Mining]\label{def:topk_problem}
Given transaction sequence $D$, base window $W_0$, base threshold $\theta_0$,
and integer $k \geq 1$, find the set $\mathcal{T}_k^* \subseteq 2^{\mathcal{I}}$
of $k$ patterns with the highest Multi-Scale Dense Coverage Scores:
\begin{equation}
\mathcal{T}_k^* = \arg\max_{\mathcal{T} \subseteq 2^{\mathcal{I}},\, |\mathcal{T}| = k}
  \min_{P \in \mathcal{T}} \mathrm{MSDCS}(P).
\end{equation}
\end{definition}

% ============================================================
\section{Theoretical Results}\label{sec:theorems}
% ============================================================

\begin{theorem}[Anti-Monotonicity of Dense Coverage]\label{thm:antimonotone}
For any patterns $C \subseteq P$ and any window size $W$ and threshold $\theta$:
\begin{equation}
\mathrm{DCS}(P, W, \theta) \leq \mathrm{DCS}(C, W, \theta).
\end{equation}
\end{theorem}

\begin{proof}
Since $C \subseteq P$, we have $P \subseteq d_t \Rightarrow C \subseteq d_t$,
hence $s_P^W(t) \leq s_C^W(t)$ for all $t$.

\textbf{Step 1}: Every dense interval $[a,b]$ of $P$ is contained in
some dense interval of $C$. If $s_P^W(t) \geq \theta$ for $t \in [a,b]$,
then $s_C^W(t) \geq s_P^W(t) \geq \theta$, so $[a,b]$ is also dense for $C$
(possibly as part of a larger interval).

\textbf{Step 2}: Within any dense interval, the per-position support
satisfies $s_P^W(t) \leq s_C^W(t)$.

\textbf{Step 3}: Combining Steps 1 and 2:
$\mathrm{DCS}(P, W, \theta)
= \sum_{[a,b] \in \mathcal{D}(P)} \sum_{t=a}^{b} s_P^W(t)
\leq \sum_{[a,b] \in \mathcal{D}(C)} \sum_{t=a}^{b} s_C^W(t)
= \mathrm{DCS}(C, W, \theta)$.
\end{proof}

\begin{corollary}[Multi-Scale Anti-Monotonicity]\label{cor:ms_antimonotone}
For $C \subseteq P$: $\mathrm{MSDCS}(P) \leq \mathrm{MSDCS}(C)$.
\end{corollary}

\begin{proof}
Direct consequence of Theorem~\ref{thm:antimonotone} applied at each scale level,
combined with non-negative weights $\omega_\ell \geq 0$.
\end{proof}

\begin{theorem}[B\&B Correctness]\label{thm:bb_correct}
Algorithm~\ref{alg:topk_bb} returns the exact top-$k$ patterns
with highest MSDCS scores.
\end{theorem}

\begin{proof}
The algorithm maintains a min-heap $H$ of size $k$ with the current
$k$-th best score $\tau$.
At each node in the itemset lattice, it computes
$\overline{\mathrm{MSDCS}}(P) = \mathrm{MSDCS}(C)$ where $C$ is the
current prefix subset of candidate $P$.
By Corollary~\ref{cor:ms_antimonotone},
$\mathrm{MSDCS}(P) \leq \overline{\mathrm{MSDCS}}(P) = \mathrm{MSDCS}(C)$.
If $\mathrm{MSDCS}(C) \leq \tau$, then no extension of $C$ can enter the top-$k$,
and the subtree rooted at $C$ is safely pruned.

Since the upper bound is valid (Corollary~\ref{cor:ms_antimonotone})
and the algorithm explores all non-pruned branches exhaustively,
the result is exact.
\end{proof}

\begin{theorem}[Dyadic Decomposition Completeness]\label{thm:dyadic_complete}
For any window size $W \in [W_0, 2^L \cdot W_0]$, there exist adjacent
dyadic levels $\ell, \ell+1$ with $W_\ell \leq W \leq W_{\ell+1}$
and $W_{\ell+1} / W_\ell = 2$.
A pattern that is dense at scale $W$ with density ratio
$\theta / W$ is necessarily dense at either $W_\ell$ or $W_{\ell+1}$
(or both) with the corresponding proportional threshold.
\end{theorem}

\begin{proof}
Let $W \in [2^\ell W_0, 2^{\ell+1} W_0)$ for some $\ell$.
Suppose $s_P^W(t) \geq \theta$ where $\theta/W = \rho$ is the density ratio.

Consider the window $[t, t+W)$. It can be covered by at most two
adjacent windows of size $W_{\ell+1} = 2^{\ell+1} W_0 \geq W$:
the window $[t, t+W_{\ell+1})$ contains $[t, t+W)$, so
$s_P^{W_{\ell+1}}(t) \geq s_P^W(t) \geq \theta$.
Since $\theta_{\ell+1} = \lceil \rho \cdot W_{\ell+1} \rceil
\geq \lceil \rho \cdot W \rceil = \theta$,
and $s_P^{W_{\ell+1}}(t) \geq \theta \geq \theta$ (the count can
only increase with a larger window), the position $t$ is dense at level $\ell+1$
if $s_P^{W_{\ell+1}}(t) \geq \theta_{\ell+1}$.

More precisely, the additional transactions in
$[t+W, t+W_{\ell+1})$ contribute non-negative support.
Thus $s_P^{W_{\ell+1}}(t) \geq s_P^W(t)$ and the density is captured
at the coarser scale with at most a factor-of-2 resolution loss.
\end{proof}

% ============================================================
\section{Complexity Analysis}\label{sec:complexity}
% ============================================================

\begin{theorem}[Time Complexity]\label{thm:complexity}
Let $N$ be the number of transactions, $m$ the number of items,
$L = O(\log(N/W_0))$ the number of dyadic levels,
and $F$ the number of frequent singletons.

\textbf{Worst case}:
The exhaustive search explores $O(2^F)$ candidate patterns,
each requiring $O(L \cdot N)$ work for multi-scale dense interval computation.
Total: $O(2^F \cdot L \cdot N)$.

\textbf{With B\&B pruning}:
Let $\mathcal{P}$ be the set of patterns actually evaluated
(not pruned). Then the complexity is $O(|\mathcal{P}| \cdot L \cdot N)$.
Empirically, $|\mathcal{P}| \ll 2^F$ due to early termination.

\textbf{Per-scale dense interval computation}:
Using the sliding window algorithm from Phase~1,
computing $\mathcal{D}(P, W, \theta)$ takes $O(|T_P|)$ time
where $T_P$ is the set of timestamps where $P$ occurs.
\end{theorem}

\begin{theorem}[Space Complexity]\label{thm:space}
The algorithm requires:
\begin{itemize}
\item $O(N)$ for the transaction database,
\item $O(m \cdot \bar{f})$ for item-to-timestamp maps
  (where $\bar{f}$ is average item frequency),
\item $O(k)$ for the top-$k$ heap,
\item $O(L \cdot \max_P |\mathcal{D}(P)|)$ for intermediate
  dense intervals across scales.
\end{itemize}
Total: $O(N + m \cdot \bar{f} + k + L \cdot I_{\max})$
where $I_{\max}$ is the maximum number of intervals for any pattern.
\end{theorem}
